% Track e3 -- Up the ladder: the rung Cl(3n,3n).
% Verification: investigations/emergence/verify/e3_ladder.py (claims e3-1a ... e3-6c).
% New references: sections/bib_e3.tex (Howe89). Also cited: Chevalley54, LawsonMichelsohn89, Pauli40, WeinbergQFT (bib_e1.tex),
% EhrenfestOppenheimer31 (bib_e2.tex).
\section{Up the ladder}
\label{e3:sec:main}

\Cref{e1:sec:main,e2:sec:main} treat the second rung. This section treats the rung $\Cl(3n,3n)$ for every $n$ and the step from $n$ to
$n+1$. The question is the same as at $n=2$: in what precise sense, and under which added structures, does rung $n$ carry $n$ particles
rather than one particle with a larger internal space, and what changes from one rung to the next. The answers, with their classes, are
these. (1) $\Cl(3n,3n)=M_{8^n}(\R)$ at every rung. The step embeds rung $n$ in rung $n+1$ with multiplicity eight and adds three Fock
modes, one frame and one copy direction; the commuting copy of the new constituent is its generators dressed by the volume element of
rung $n$ (fixed; \cref{e3:sec:rung}). (2) For $n$ equal masses the graded rest Hamiltonian is one clock of mass $\sqrt n\,m$, and the
ungraded one is $n$ clocks of mass $m$, with spectrum $m(n-2j)$ of multiplicity $\binom nj4^n$ (fixed, partly restated from
\cref{e1:prop:rest}; \cref{e3:sec:mass}). (3) A rotation through $2\pi$ of the diagonal interior Lorentz action is $(-1)^n$ on $S_n$, and one
of the diagonal field-strength action is $+\id$. The interior composite is a fermion exactly at odd $n$, and the field-strength composite
has integer spin at every rung (restated from Proposition~10.9 of \cite{ES17}, with the second half fixed; \cref{e3:sec:spin}). (4) On vectors
the copy rotations of all six directions span the commutant of the diagonal $\so{3,3}$, and the copy group $O(n)$ has the diagonal $\so{3,3}$
as its commutant. On $S_n$ the Casimirs satisfy $C_{\rm cop}+C_{\rm diag}=\tfrac34n(n+4)$. This is Howe's skew duality $(O(n),\so6)$ on the spinor module of $\so{6n}$, and at
$n=2,3$ the joint decomposition has the predicted multiplicity-free form (fixed; the duality inherited; \cref{e3:sec:dual}). (5) The Koszul
swaps realise the permutation matrices of that $O(n)$, twisted by the sign character. Their symmetric and antisymmetric sectors are mirror images, of equal dimension and
never empty. The plain permutations of the common complex space give $\Lambda^n_\C$ of the Dirac module, which vanishes for $n\ge5$ (fixed;
the exterior powers inherited; \cref{e3:sec:sym}). (6) The algebra singles out no rung, and its invariant forms single out only the parity of $n$, through a period of four. Odd rungs are
singled out by the interior Lorentz action, even rungs by the binding action, and $n=3$ by the projector $P_*$ (restated and fixed; \cref{e3:sec:table}). (7) Going up the ladder
enlarges one particle; it becomes the addition of a particle only through the constructed structures of \cref{e1:sec:main}, now needed at every
rung (interpretive, drawn from the fixed statements; \cref{e3:sec:answer}). Identifiers e3-$n$ refer to the \texttt{PASS} lines of
\texttt{verify/e3\_ladder.py} (\cref{e3:tab:claims}).

\paragraph{Notation.} As in \cref{e1:sec:main}: $S=\R^8$, $S_n=S^{\otimes n}$, $E_{(k,a)}=\omega^{\otimes(k-1)}\otimes e_a\otimes\id^{\otimes(n-k)}$, and
$V_n=\R^n\otimes V=\R^{3n,3n}$ with $\eta_n=\id_n\otimes\eta$ and $G_n=\id_n\otimes G$. $\Gamma^{(k)}=E_{(k,1)}\cdots E_{(k,6)}=\id\otimes\cdots\otimes\omega\otimes\cdots\otimes\id$ is the
volume element of constituent $k$, and $\Gamma_n=\Gamma^{(1)}\cdots\Gamma^{(n)}$. The ungraded generators are
$u_{(k,a)}=\id^{\otimes(k-1)}\otimes e_a\otimes\id^{\otimes(n-k)}$, $J^{(k)}=E_{(k,2)}E_{(k,3)}$, and $a^\pm_{k,i}=\half(E_{(k,i)}\pm E_{(k,i+3)})$, $N=\sum_{k,i}a^+_{k,i}a^-_{k,i}$
(Proposition~7.2 of \cite{RE}); $|0\rangle$ is the first basis vector of $S_n$, the Fock vacuum. The diagonal generators are
$\Delta_{ab}=\half\sum_kE_{(k,a)}E_{(k,b)}$. All fifteen span the \emph{diagonal} $\so{3,3}$; those with $a,b\in\{1,4,5,6\}$ span the diagonal $\lint$; and the
paired rotations $\Delta_{ij}-\Delta_{i+3,j+3}$ and paired boosts $\Delta_{i,j+3}-\Delta_{j,i+3}$ span the diagonal $\lG$ (Theorem~12.1 of \cite{ES17}). The
copy rotations are $C^{kl}=\half\sum_a\eta_{aa}E_{(k,a)}E_{(l,a)}$ for $k<l$, with $C^{12}=C_b$ of Proposition~7.3(ii) of \cite{RE}, and $M_{kl}$ are their
interior parts (Section~10.1 of \cite{ES17}). The Casimirs are $C_{\rm diag}=-\sum_{a<b}\eta^{aa}\eta^{bb}\Delta_{ab}^2$ over all six directions,
$C_\Delta$ the same over the interior four (Section~10.2 of \cite{ES17}), $C_G=-\sum_iL_i^2+\sum_iK_i^2$ for the paired rotations $L_i$ and boosts
$K_i$ of $\lG$, $C_{\rm cop}=-\sum_{k<l}(C^{kl})^2$ and $C^{\rm int}_{\rm cop}=-\sum_{k<l}M_{kl}^2$. A Lorentz Casimir takes the value
$2[j_L(j_L+1)+j_R(j_R+1)]$ on $(j_L,j_R)$. We write $\mathfrak S_n$ for the symmetric group, to keep $S_n$ for the module.

\subsection{The rung and the step}
\label{e3:sec:rung}

\begin{proposition}[The rung and the step]
\label{e3:prop:rung}
\begin{enumerate}[label=(\roman*),leftmargin=2.2em]
\item Every blade of $\Cl(3n,3n)$ other than $\pm\id$ is traceless on $S_n$. Hence the $64^n$ blades are linearly independent and
$\Cl(3n,3n)\cong M_{8^n}(\R)$, with $S_n$ its unique irreducible module. The volume element is $E_{(1,1)}\cdots E_{(n,6)}=\Gamma_n=\omega^{\otimes n}$, and
$(-1)^N=(-1)^n\Gamma_n$.
\item $E^{(n+1)}_{(k,a)}=E^{(n)}_{(k,a)}\otimes\id$ for $k\le n$ and $E_{(n+1,a)}=\Gamma_n\otimes e_a$. So $\Cl(3n,3n)\otimes\id\subset\Cl(3n+3,3n+3)$, and
$S_{n+1}$ restricted to $\Cl(3n,3n)$ is eight copies of $S_n$. The commutant of $\Cl(3n,3n)\otimes\id$ is $\id\otimes M_8(\R)$, generated by the
ungraded generators $\id\otimes e_a=(\Gamma_n\otimes\id)E_{(n+1,a)}$ of the new constituent.
\item $N_{n+1}=N_n\otimes\id+\id\otimes N_1$, $\spec N_n=\{0,\dots,3n\}$ with multiplicities $\binom{3n}j$, and $[N,E_{(k,i)}]=E_{(k,i+3)}$,
$[N,E_{(k,i+3)}]=E_{(k,i)}$. Each step adds three modes.
\item Every diagonal bivector satisfies $\sum_{k\le n+1}E_{(k,a)}E_{(k,b)}=\bigl(\sum_{k\le n}E_{(k,a)}E_{(k,b)}\bigr)\otimes\id+\id\otimes e_ae_b$. Hence
$S_{n+1}\cong S_n\otimes S$ and $S_n\cong S^{\otimes n}$ as modules of the diagonal $\so{3,3}$, and so of the diagonal $\lint$ and $\lG$.
\end{enumerate}
\end{proposition}

\begin{proof}
(i) Moving the generators of each constituent to their own factor gives $e_I=\pm\bigotimes_je_{I_j}\omega^{p_j}$ with $p_j=|I_{j+1}|+\cdots+|I_n|$ mod $2$, and
the trace is multiplicative. On $S$, $\tr(e_J\omega^p)\ne0$ only for $(J,p)=(\emptyset,0)$ and for $J$ the full index set with $p=1$. The last
factor has $p_n=0$, so a nonzero trace needs $I_n=\emptyset$; then $p_{n-1}=0$, and so on, down to $I=\emptyset$. Trace orthogonality gives
independence, and $64^n=(8^n)^2$. The volume element is $\prod_k\Gamma^{(k)}$ because $\prod_aE_{(k,a)}=(\omega^{\otimes(k-1)})^6\otimes\omega\otimes\id$. On one
constituent $(-1)^N=-\omega$ (Proposition~7.2 of \cite{RE}), and $N=\sum_kN^{(k)}$. (ii) The formulas are equation~(2) of \cite{ES17} at rungs $n$
and $n+1$. The commutant of $M_{8^n}(\R)\otimes\id$ in $M_{8^{n+1}}(\R)$ is $\id\otimes M_8(\R)$, and $\Gamma_n^2=\id$. (iii) $N=\sum_{k,i}\half(\id-E_{(k,i)}E_{(k,i+3)})$,
the terms are commuting diagonal involutions with traceless products, and the commutators are those of Proposition~7.2 of \cite{RE}. (iv)
$E_{(n+1,a)}E_{(n+1,b)}=\Gamma_n^2\otimes e_ae_b=\id\otimes e_ae_b$.
\end{proof}

Part (i) is fixed in the realisation (e3-1a, exhaustive over all $64^n$ blades at $n\le4$). That $\Cl(p,q)\cong M_{2^{(p+q)/2}}(\R)$ for $p-q\equiv0\pmod8$
is inherited \cite{LawsonMichelsohn89}: the real structure of the algebra is the same at every rung, so it singles out none. Parts (ii) and
(iii) restate Propositions~7.1 and~7.2 of \cite{RE} and \cref{e1:prop:copies}(ii),(v) at the last factor (e3-1b, with the commutant
exhaustive over blades at $n=1,2$; e3-1c, spectrum also at $n=5$). Part (iv) is fixed (e3-6c, $n\le3$). Under (iv) every statement about the
diagonal action is a statement about tensor powers of the one-constituent module, and the step is tensoring with $S$.

\subsection{Mass up the ladder: one clock or \texorpdfstring{$n$}{n} clocks}
\label{e3:sec:mass}

Put $H_g=m\sum_kE_{(k,1)}$, built from the vectors of the rung, and $H_u=m\sum_ku_{(k,1)}$, built from the commuting copies.

\begin{proposition}[Mass up the ladder]
\label{e3:prop:mass}
\begin{enumerate}[label=(\roman*),leftmargin=2.2em]
\item The $E_{(k,1)}$ are pairwise anticommuting traceless involutions. So $H_g^2=nm^2\id$ and $\spec H_g=\{\pm\sqrt n\,m\}$, each of multiplicity
$8^n/2$. Every $J^{(k)}$ commutes with $H_g$, and $Z_g=J^{(1)}H_g$ satisfies $Z_g^2=-nm^2\id$: the rest flow $e^{-\lambda Z_g}$ is one clock of
angular frequency $\sqrt n\,m$.
\item The ungraded generators are $u_{(k,a)}=\Gamma^{(1)}\cdots\Gamma^{(k-1)}E_{(k,a)}$, with no further sign. The $u_{(k,1)}$ commute, square to $\id$, and
every product over a nonempty subset is traceless. So $\spec H_u=\{m(n-2j):j=0,\dots,n\}$, the value $m(n-2j)$ with multiplicity $\binom nj4^n$.
\item $T^{(k)}_u=J^{(k)}u_{(k,1)}=\id\otimes\cdots\otimes T\otimes\cdots\otimes\id$ are $n$ commuting complex structures: $n$ clocks of angular frequency $m$.
\item $H_g^{(n+1)}=H_g^{(n)}\otimes\id+\Gamma_n\otimes me_1$ with anticommuting summands, and $H_u^{(n+1)}=H_u^{(n)}\otimes\id+\id\otimes me_1$ with commuting
summands. The graded mass grows in quadrature, $\sqrt n\,m\to\sqrt{n+1}\,m$, and the ungraded energies add.
\end{enumerate}
\end{proposition}

\begin{proof}
(i) Graded generators of distinct constituents anticommute, and each $E_{(k,1)}$ is traceless because $\tr e_1=0$. $J^{(k)}$ is a bivector sharing no
index with any $E_{(l,1)}$, so it commutes with each (equation~(3) of \cite{ES17}), and $(J^{(k)})^2=-\id$. (ii) $\Gamma^{(1)}\cdots\Gamma^{(k-1)}=\omega^{\otimes(k-1)}\otimes\id$ and $\omega^2=\id$, so the
product is $\id^{\otimes(k-1)}\otimes e_a\otimes\id$. Products of the $u_{(k,1)}$ are tensor products of factors $e_1$ and $\id$, traceless unless all
factors are $\id$. Commuting involutions whose nonempty products are traceless have all $2^n$ joint signs with equal multiplicity $8^n/2^n=4^n$, and $j$ minus signs give
$m(n-2j)$. (iii) $J\,e_1=e_2e_3e_1=T$ on the $k$th factor. (iv) \Cref{e3:prop:rung}(ii); $\Gamma_n$ anticommutes with every generator of $\Cl(3n,3n)$.
\end{proof}

At general masses (i) and (ii) are \cref{e1:prop:rest}(vi), restated; the multiplicities, the clocks (iii) and the step (iv) are fixed
(e3-2a and e3-2b, exact over signed permutations at $n\le5$, eigenvalue counts in floating point at $n\le3$; e3-6c). At $n=2$, (i) is the
Dirac duplication of Proposition~8.11(iv) of \cite{RE}. So rung $n$ read through its vectors carries one Dirac mass, $\sqrt n\,m$, along the unit
vector $n^{-1/2}\sum_kE_{(k,1)}$ of $\R^{3n,3n}$, with one clock; read through its commuting copies it carries $n$ masses $m$ with $n$ clocks and
threshold $nm$. The quadrature is the statement that the copies of the graded tower are directions of one momentum space (\cref{e1:sec:rest}).

\subsection{Spin parity}
\label{e3:sec:spin}

Let $R_n=\half\sum_kE_{(k,4)}E_{(k,5)}$, the rotation about the third axis of the diagonal $\lint$, and
$L_n=\half\sum_k\bigl(E_{(k,1)}E_{(k,2)}-E_{(k,4)}E_{(k,5)}\bigr)$, the paired rotation of the $(1,2)$ and $(4,5)$ planes in the diagonal $\lG$. Both
generate rotations of period $2\pi$ on vectors.

\begin{proposition}[Spin parity up the ladder]
\label{e3:prop:spin}
\begin{enumerate}[label=(\roman*),leftmargin=2.2em]
\item $R_n$ is a sum of $n$ commuting terms $\half b_k$ with $b_k^2=-\id$. Hence $e^{2\pi R_n}=(-1)^n\id$, and the eigenvalues of $2R_n$ are $i$ times
integers congruent to $n$ modulo $2$. Under the diagonal $\lint$ the whole of $S_n$ has half-integer spin for odd $n$ and integer spin for even
$n$. The Casimir spectra are
\[
\begin{array}{@{}ll@{}}
n=1: & C_\Delta=\tfrac32\,(8),\\
n=2: & 0\,(8),\ 3\,(32),\ 4\,(24),\\
n=3: & \tfrac32\,(160),\ \tfrac{11}2\,(288),\ \tfrac{15}2\,(64),\\
n=4: & 0\,(160),\ 3\,(1024),\ 4\,(864),\ 8\,(864),\ 9\,(1024),\ 12\,(160),
\end{array}
\]
where the values $0,\tfrac32,3,4,\tfrac{11}2,8,\tfrac{15}2,9,12$ are those of $(0,0)$, $(\half,0)\oplus(0,\half)$, $(\half,\half)$, $(1,0)\oplus(0,1)$, $(1,\half)\oplus(\half,1)$, $(1,1)$, $(\tfrac32,0)\oplus(0,\tfrac32)$,
$(\tfrac32,\half)\oplus(\half,\tfrac32)$ and $(2,0)\oplus(0,2)$.
\item $L_n$ commutes with $G_n$ on vectors and is a sum of $2n$ commuting terms $\pm\half b$ with $b^2=-\id$. Hence $e^{2\pi L_n}=+\id$ at every rung. The
diagonal $\lG$ has Casimir $C_G=3$ on all of $S$; $0,4,8$ on $4,24,36$ dimensions at $n=2$; and $3,9,15$ on $128,256,128$ at $n=3$: the content of
$(\half,\half)^{\otimes n}$, integer spin at every rung.
\item The central element $-\id$ of $\mathrm{Spin}(3,3)$, which is $e^{\pi e_4e_5}=e^{\pi e_1e_2}$ on $S$, acts on $S_n$ through the diagonal
$\mathrm{Spin}(3,3)$ as $(-1)^n$. The rotation through $2\pi$ of $\lint$ is this element, and that of $\lG$ is not.
\item For the actions of Theorem~10.8 of \cite{ES17}, with $\lint$ on $k_I$ constituents, $\lG$ on $k_B$ and trivially on the rest, the rotation
through $2\pi$ is $(-1)^{k_I}$.
\end{enumerate}
\end{proposition}

\begin{proof}
(i) $(E_{(k,4)}E_{(k,5)})^2=-E_{(k,4)}^2E_{(k,5)}^2=-\id$, and bivectors of distinct constituents share no index and commute. So
$e^{2\pi R_n}=\prod_k(\cos\pi+b_k\sin\pi)$. Each listed value of $2[j_L(j_L+1)+j_R(j_R+1)]$ is attained by exactly one unordered pair $(j_L,j_R)$, which identifies the
modules. (ii) $e_ie_j-e_{i+3}e_{j+3}$ commutes with $G$ (Proposition~7.8(i) of \cite{RE}). The $2n$ terms share no index, and
$e^{2\pi L_n}=\prod_k(-\id)(-\id)$. By \cref{e3:prop:rung}(iv), $S_n\otimes\C\cong(S\otimes\C)^{\otimes n}=2^n(\half,\half)^{\otimes n}$ under the diagonal $\lG$,
since $S\otimes\C=2(\half,\half)$ (Proposition~13.12 of \cite{ES17}); the Clebsch--Gordan series gives the listed values. (iii) The diagonal
$\mathrm{Spin}(3,3)$ acts on $S^{\otimes n}$ by $g\mapsto g^{\otimes n}$ (\cref{e3:prop:rung}(iv)). $e^{2\pi\cdot\frac12e_4e_5}$ is $-\id$, while the double
rotation $e^{2\pi\cdot\frac12(e_1e_2-e_4e_5)}=e^{\pi e_1e_2}e^{-\pi e_4e_5}$ is $+\id$. (iv) Products of the factors of (i) and (ii) over the respective
constituents.
\end{proof}

Part (i) restates Proposition~10.9 and Proposition~10.10(i) of \cite{ES17} (fixed there), with the element and the spectra re-verified (e3-3a,
exact at $n\le5$; e3-3c, floating-point spectra at $n\le4$). Part (ii) extends Proposition~13.12 of \cite{ES17} from $n=1$ to every rung and is
fixed (e3-3b, e3-3c). Part (iii) is fixed (e3-3b), and (iv) restates Proposition~10.9 of \cite{ES17}. The parity of (i) is the
Ehrenfest--Oppenheimer rule, a composite of $n$ spin-$\half$ particles is a fermion iff $n$ is odd (inherited \cite{EhrenfestOppenheimer31}).
The rule holds equally for one particle whose internal space is the $n$-fold tensor product of spinors. So the spin parity does not decide
between one particle and $n$. Under the field-strength reading it does not arise: every rung carries integer spin, $\mathrm{Spin}(3,3)\cong SL(4,\R)$
restricted to $\lG$ integrates to $SO^+(1,3)$ on every $S_n$, and the diagonal $\lG$ composite is a boson at every rung.

\begin{remark}[Comparison with Theorem~13.15 of \cite{ES17}]
\label{e3:rem:copyparity}
Theorem~13.15 of \cite{ES17} says that a copy-invariant four-vector in $\Lambda^g(\R^n\otimes\R^{1,3})$ needs $n$ odd and $g$ odd, $n\le g\le3n$. The
statement concerns the algebra and (i) concerns the module, and they are logically independent. Both are products over the $n$ constituents:
in (i) each constituent contributes a factor $-1$ to the $2\pi$ rotation, and in the theorem the copy invariant of odd grade needs the
Levi-Civita symbol of $SO(n)$ with its $n$ indices. They meet on the Dirac block. At odd $n$ the module contains spin $\half$, and at $n=3$ the
copy-singlet four-vector $Q_\mu$ of grade three builds a gauge-invariant $\Cl(1,3)$ on that block (Theorem~10.6 of \cite{ES17}). The check e3-6b
re-verifies that $Q_\mu$ commutes with the interior copy rotations $M_{kl}$ and satisfies $[\Delta_{ab},Q_c]=\eta_{bc}Q_a-\eta_{ac}Q_b$, and that the $n=2$
analogue $\sum_\nu\eta_{\nu\nu}(\gamma^{(1)}_\mu\gamma^{(2)}_\nu+\gamma^{(1)}_\nu\gamma^{(2)}_\mu)$ commutes with $M_{12}$ but is not a vector. At even $n$ there is
neither a fermion to carry gamma matrices nor a copy-invariant four-vector to build them. Under the diagonal $\lG$ neither statement has
an analogue: the spin is integral at every rung, and grade one is $n[(1,0)\oplus(0,1)]$, $n$ field strengths and no four-vector (Section~13.5
of \cite{ES17}). The juxtaposition is interpretive; the facts are restated and fixed.
\end{remark}

\subsection{A dual pair up the ladder}
\label{e3:sec:dual}

\begin{theorem}[The copy algebra and the diagonal $\so{3,3}$]
\label{e3:thm:dual}
\begin{enumerate}[label=(\roman*),leftmargin=2.2em]
\item On vectors, the commutant of the diagonal $\so{3,3}$ in $\so{3n,3n}$ is $\so n\otimes\id_6$, of dimension $\binom n2$, spanned by the vector
actions of the $C^{kl}$. The commutant of $\so n\otimes\id_6$ together with the copy permutations is $\id_n\otimes\so{3,3}$ for $n\ge2$; that of
$\so n\otimes\id_6$ alone has dimension $36$ at $n=2$ and $15$ for $n\ge3$. The weights $\eta_{aa}$ are forced: $\sum_aE_{(k,a)}E_{(l,a)}$ does not commute
with the diagonal $\so{3,3}$.
\item On $S_n$, $C^{kl}$ is the spin lift of $(e_ke_l^{\!\top}-e_le_k^{\!\top})\otimes\id_6$: $[C^{kl},E_{(l,a)}]=E_{(k,a)}$ and $[C^{kl},E_{(k,a)}]=-E_{(l,a)}$. The
$C^{kl}$ commute with the fifteen $\Delta_{ab}$ and close on $\so n$, with $[C^{kl},C^{lm}]=C^{km}$. Each is a sum of six commuting terms of square
$-\tfrac14$, so $e^{2\pi C^{kl}}=\id$. In Fock form
\[
C^{kl}=\sum_{i=1}^3\bigl(a^+_{k,i}a^-_{l,i}-a^+_{l,i}a^-_{k,i}\bigr),\qquad C^{kl}a^+_{l,i}|0\rangle=a^+_{k,i}|0\rangle,
\]
so the copy algebra conserves $N$ and moves quanta between constituents.
\item As operators on $S_n$,
\[
C_{\rm cop}+C_{\rm diag}=\tfrac34\,n(n+4)\,\id,
\]
that is $\tfrac{15}4,9,\tfrac{63}4,24,\tfrac{135}4,\tfrac{81}2$ at $n=1,\dots,6$. The interior identity of Theorem~10.4 of \cite{ES17},
$C^{\rm int}_{\rm cop}+C_\Delta=\half n(n+2)\,\id$, holds at every $n$. Both constants are $nd(d+n-2)/8$, at $d=6$ and $d=4$.
\item At $n=2$ the commutant of the diagonal $\so{3,3}$ in $M_{64}(\R)$ has dimension $14$ and is spanned by the words in $C^{12}$ and the Koszul
swap $\tau_K$, with $C^{12}\tau_K=-\tau_KC^{12}$: it is the image of the group algebra of $O(2)$. The words in $C^{12}$ alone span $7$ dimensions.
\end{enumerate}
\end{theorem}

\begin{proof}
(i) $V=\R^{3,3}$ is absolutely irreducible under $\so{3,3}$, so by Schur's lemma the commutant of $\id_n\otimes\so{3,3}$ in $\End(\R^n\otimes V)$ is
$\End(\R^n)\otimes\id_6$. An element $A\otimes\id_6$ is $\eta_n$-antisymmetric iff $A^{\!\top}=-A$. Conversely, $\R^n$ is absolutely irreducible under
the group generated by $SO(n)$ and the permutation matrices for $n\ge2$, and under $\so n$ alone for $n\ge3$; at $n=2$ the commutant of $\so2$
is $\operatorname{span}\{\id,\varepsilon\}$, which adds $\varepsilon\otimes\mathrm{Sym}_\eta(V)$, of dimension $21$. The unweighted sum lifts
$(e_ke_l^{\!\top}-e_le_k^{\!\top})\otimes\eta$, and $\eta$ does not commute with the boosts. (ii) The spin lift $\sigma(X)=\tfrac14\sum_{A,B}X_{AB}\eta_{BB}E_AE_B$
satisfies $[\sigma(X),E_C]=\sum_AX_{AC}E_A$ and is a Lie algebra homomorphism; the vector matrices commute by (i). The six terms
$\half\eta_{aa}E_{(k,a)}E_{(l,a)}$ share no index, and $(E_{(k,a)}E_{(l,a)})^2=-\eta_{aa}^2\id$. For the Fock form substitute
$E_{(k,i)}=a^+_{k,i}+a^-_{k,i}$, $E_{(k,i+3)}=a^+_{k,i}-a^-_{k,i}$: $C^{kl}=\sum_i(a^+_{k,i}a^-_{l,i}+a^-_{k,i}a^+_{l,i})$, and $a^-_{k,i}a^+_{l,i}=-a^+_{l,i}a^-_{k,i}$ for
$k\ne l$ by the CAR of Proposition~7.2 of \cite{RE}. (iii) The proof of Theorem~10.4 of \cite{ES17} uses neither $d=4$ nor the signature of the
interior: with $B^{kl}_{ab}=E_{(k,a)}E_{(k,b)}E_{(l,a)}E_{(l,b)}$ it gives $C_{\rm diag}=\tfrac{nd(d-1)}8-\half\mathsf S$ and $C_{\rm cop}=\binom n2\tfrac d4+\half\mathsf S$
with the same cross term $\mathsf S=\sum_{k<l}\sum_{a<b}\eta^{aa}\eta^{bb}B^{kl}_{ab}$. At $d=6$ the constant is $\tfrac{15n}4+\tfrac{3n(n-1)}2=\tfrac34n(n+4)$.
(iv) Computed: the nullity of $\sum_X\ad_X^{\!\top}\ad_X$ over the fifteen generators, and the rank of the words.
\end{proof}

Part (i) is fixed (e3-4a, exact over $\Q$ at $n\le5$); at $n=3$ its interior analogue is Theorem~10.2 of \cite{ES17}. Part (ii) restates
Proposition~7.3(ii) of \cite{RE} at $n=2$ and Proposition~9.6(i) of \cite{RE}, and is fixed at $n\le4$ (e3-4b, e3-4f); the Fock form at $n=2$ is
\cref{e2:prop:hop}(i) with the opposite orientation. Part (iii) is fixed: the constant $\tfrac34n(n+4)$ is new, and the identity was checked exactly
as matrices on $S_n$ at $n\le4$ and blade by blade in the abstract Clifford algebra at $n\le6$; the interior identity is restated and checked in
the same two ways (e3-4c). Part (iv) is fixed (e3-4e; nullity in floating point with gap $>1$, rank exact). The $\eta$-weighting is where the
signature enters: a copy rotation pairs two generators of the same square, so it is compact (Theorem~10.2 of \cite{ES17}), and $C_{\rm cop}\ge0$.

\paragraph{Howe's skew duality.} By Proposition~7.2 of \cite{RE}, $S_n\otimes\C$ is the exterior algebra on the $3n$ creation operators
$a^+_{k,i}$, that is $\Lambda(\C^n\otimes\C^3)$ with $k$ the copy index and $i$ the mode index. The complexified diagonal $\so{3,3}$ is spanned by
the quadratics $\sum_kx_ky_k$ with $x,y\in\{a^\pm_{\cdot,i}\}$, the $O(n)$-invariant contractions over the copy index, which form, up to multiples
of $\id$ and by the substitution used in (ii),
$\Lambda^2\C^3\oplus\mathfrak{gl}(3)\oplus\Lambda^2\C^3\cong\so{6,\C}$. By (ii) the copy rotations are the $\so n\subset\mathfrak{gl}(n)$ acting on the copy index.
This is the setting of Howe's skew $(O(n),\so{2m})$ duality on $\Lambda(\C^n\otimes\C^m)$, here with $m=3$ \cite{Howe89}: the two actions generate mutual
commutants, and the module decomposes without multiplicity as $\bigoplus_\lambda\rho_\lambda\otimes\pi_{\lambda'}$, with $\rho_\lambda$ an irreducible
representation of $O(n)$ and $\pi_{\lambda'}$ one of $\so6$ determined by it. The theorem, and the linear relation between the two Casimirs, are
inherited; (iii) is its constant and (iv) its commutant at $n=2$ in the tower. The $O(n)$ contains more than $SO(n)$: at $n=2$ the reflection is
$\tau_K$, and \cref{e3:sec:sym} shows that the Koszul swaps realise the permutation matrices at every rung.

\begin{table}[t]
\centering
\footnotesize
\begin{tabular}{@{}cccclcc@{}}
\toprule
$n$ & copy label & $C_{\rm cop}$ & $C_{\rm diag}$ & $\so6$ modules $(\lambda_1,\lambda_2,\lambda_3)$ & dimensions & rank in $S_n$\\
\midrule
$2$ & $|j|=0$ & $0$ & $9$ & $(1,1,1)$, $(1,1,-1)$ & $10+10$ & $1\cdot20=20$\\
 & $|j|=1$ & $1$ & $8$ & $(1,1,0)$ & $15$ & $2\cdot15=30$\\
 & $|j|=2$ & $4$ & $5$ & $(1,0,0)$ & $6$ & $2\cdot6=12$\\
 & $|j|=3$ & $9$ & $0$ & $(0,0,0)$ & $1$ & $2\cdot1=2$\\
\midrule
$3$ & $j=0$ & $0$ & $\tfrac{63}4$ & $(\tfrac32,\tfrac32,\pm\tfrac32)$ & $20+20$ & $1\cdot40=40$\\
 & $j=1$ & $2$ & $\tfrac{55}4$ & $(\tfrac32,\tfrac32,\pm\tfrac12)$ & $36+36$ & $3\cdot72=216$\\
 & $j=2$ & $6$ & $\tfrac{39}4$ & $(\tfrac32,\tfrac12,\pm\tfrac12)$ & $20+20$ & $5\cdot40=200$\\
 & $j=3$ & $12$ & $\tfrac{15}4$ & $(\tfrac12,\tfrac12,\pm\tfrac12)$ & $4+4$ & $7\cdot8=56$\\
\bottomrule
\end{tabular}
\caption{Joint decomposition of $S_n$ under the copy algebra and the diagonal $\so{3,3}$ (\cref{e3:thm:dual}). At $n=2$, $j$ is the
$\so2$ weight; at $n=3$, $j$ is the copy isospin. $C_{\rm diag}=\langle\lambda,\lambda+2\rho\rangle$ with $\rho=(2,1,0)$, so that the spinor has
$\tfrac{15}4$ and the vector $5$. Ranks are real dimensions of $S_n$, equal to complex dimensions of $S_n\otimes\C$. Check e3-4d.}
\label{e3:tab:joint}
\end{table}

\begin{proposition}[Joint decomposition]
\label{e3:prop:joint}
The joint spectrum of $(C_{\rm cop},C_{\rm diag})$ on $S_2$ and $S_3$ is that of \cref{e3:tab:joint}. Every rank is the dimension of the copy module
times the sum of the dimensions of the listed $\so6$ modules, and every $\so6$ Casimir equals $\tfrac34n(n+4)-C_{\rm cop}$. At $n=4$,
$C_{\rm cop}=0,3,4,8,9,12,15,16,19,24$ on $70,560,504,810,1024,200,320,450,144,14$ dimensions, the values $2[j_+(j_++1)+j_-(j_-+1)]$ of the
$\so4=\su2\oplus\su2$ modules $(0,0)$, $(\half,\half)$, $(1,0)\oplus(0,1)$, $(1,1)$, $(\tfrac32,\half)\oplus(\half,\tfrac32)$, $(2,0)\oplus(0,2)$, $(\tfrac32,\tfrac32)$,
$(2,1)\oplus(1,2)$, $(\tfrac52,\half)\oplus(\half,\tfrac52)$, $(3,0)\oplus(0,3)$, each paired with $\so6$ modules of Casimir $24-C_{\rm cop}$ whose
dimensions reproduce the ranks.
\end{proposition}

\begin{proof}
At $n=2,3$ the minimal polynomial of $C_{\rm cop}$ is checked exactly and the ranks are traces of its Lagrange projectors, in integer
arithmetic. The $\so6$ dimensions are the Weyl dimension formula of $D_3$. At $n=4$, $C_{\rm cop}$ commutes with $N$ (by (ii)) and is
diagonalised on each of the thirteen number sectors.
\end{proof}

This is fixed (e3-4d), and the $n=3$ copy spectrum restates Proposition~9.6(ii) of \cite{RE}. The joint spectra lie on the line of
\cref{e3:thm:dual}(iii). Copy type and $\so{3,3}$ type determine each other, as $j_{\rm cop}$ and Lorentz type do in Corollary~10.5 of
\cite{ES17}. Two features recur up the ladder. The sector of largest copy Casimir carries the smallest $\so{3,3}$ module: the trivial module at
$n=2,4$ and the half-spinors $\mathbf4,\bar{\mathbf4}$ of one constituent at $n=3$. And the copy isospin is integral at $n=2,3$, while at $n=4$ the
modules of $\so4=\su2\oplus\su2$ include doublets of each factor, always with $j_++j_-$ integral, as single-valuedness requires (\cref{e3:thm:dual}(ii)).

\subsection{The symmetric group}
\label{e3:sec:sym}

Let $\tau_k=\id^{\otimes(k-1)}\otimes\tau_K\otimes\id^{\otimes(n-k-1)}$, with $\tau_K=D\,P$ the Koszul swap of Proposition~7.3(iii) of \cite{RE}. For a
surplus direction $s\in\{2,3\}$ let $\Theta_s=E_{(1,s)}E_{(2,s)}\cdots E_{(n,s)}$, and let $C_3=\sum_{a<b,\,c<d,\,e<f}\varepsilon_{abcdef}L^{ab}L^{cd}L^{ef}$ be the cubic invariant of
the diagonal $\so{3,3}$, with $L^{ab}=\eta^{aa}\eta^{bb}\Delta_{ab}$.

\begin{proposition}[The Koszul symmetric group]
\label{e3:prop:koszul}
\begin{enumerate}[label=(\roman*),leftmargin=2.2em]
\item $\tau_k^2=\id$, $(\tau_k\tau_{k+1})^3=\id$ and $\tau_k\tau_l=\tau_l\tau_k$ for $|k-l|>1$, so the $\tau_k$ define an action of $\mathfrak S_n$ on $S_n$. Moreover
$\tau_kE_{(j,a)}\tau_k=E_{(s_kj,a)}$ and $\tau_k=\tfrac18\prod_a(E_{(k,a)}-E_{(k+1,a)})$: $\tau_k$ is the lift to $\mathrm{Pin}(3n,3n)$ of the transposition matrix of
copies $k,k+1$, a product of six reflections. It commutes with the diagonal $\so{3,3}$, with $N$, with $H_g$ and with $C_{\rm cop}$, and conjugates
$C^{kl}$ as the permutation does. On the vacuum $\tau_k|0\rangle=-|0\rangle$.
\item At $n=3$ the $\mathfrak S_3$-isotypic parts have the content
\begin{center}
\begin{tabular}{@{}lccc|cccc@{}}
\toprule
 & $C_\Delta=\tfrac32$ & $\tfrac{11}2$ & $\tfrac{15}2$ & $C_{\rm cop}=0$ & $2$ & $6$ & $12$\\
\midrule
trivial & $24$ & $48$ & $16$ & $20$ & $36$ & $20$ & $12$\\
sign & $24$ & $48$ & $16$ & $20$ & $36$ & $20$ & $12$\\
standard & $112$ & $192$ & $32$ & $0$ & $144$ & $160$ & $32$\\
\bottomrule
\end{tabular}
\end{center}
(real dimensions; the standard rows are twice the multiplicity of the two-dimensional representation). The Dirac block, $C_\Delta=\tfrac32$,
contains $24$ totally symmetric and $24$ totally antisymmetric dimensions.
\item $\Theta_s$ is invertible, anticommutes with every $\tau_k$, commutes with the diagonal $\lint$ and with every $C^{kl}$, and anticommutes with $C_3$. It
maps the Koszul-symmetric part isomorphically onto the Koszul-antisymmetric part, as modules of $\lint\oplus\so n$, and exchanges the two
chiralities of $\so6$. At $n=3$ the eigenvalues of $8C_3$ are $-630,-210,-126$ on the symmetric part of the copy sectors $j=0,1,2$ and
$+630,+210,+126$ on the antisymmetric part.
\item The Koszul-symmetric part and the Koszul-antisymmetric part of $S_n$ have the same dimension, $\sum_k\binom{n-k+3}3\binom4k$, which is
$8,32,88,192,360$ at $n=1,\dots,5$. Neither is ever zero.
\end{enumerate}
\end{proposition}

\begin{proof}
(i) The relations are checked. $(E_{(k,a)}-E_{(k+1,a)})^2=2\eta_{aa}$, and the reflection in this vector exchanges $E_{(k,a)}$ and $E_{(k+1,a)}$; the product
over the six directions acts on vectors as the transposition matrix $P$, because the sign $(-1)^6$ of the twisted adjoint action is $+1$. Hence
$\tau_k\sigma(X)\tau_k^{-1}=\sigma(PXP^{-1})$ for every $X\in\so{\eta_n}$ and $\tau_kv\tau_k^{-1}=Pv$ for every vector. $P$ commutes with $\id_n\otimes\so{3,3}$ and with
$K_n$, and fixes $\sum_kE_{(k,1)}$, so $\tau_k$ commutes with the $\Delta_{ab}$, with $N=\tfrac{3n}2+\sigma(K_n)$ (\cref{e2:lem:N}(i)) and with $H_g$; it permutes
the $C^{kl}$ up to sign and so fixes $C_{\rm cop}$. (ii) Traces of the isotypic projectors
$\frac{\dim\rho}{6}\sum_\sigma\chi_\rho(\sigma)\tau_\sigma$ times the spectral projectors of $C_\Delta$ and $C_{\rm cop}$, in exact integer arithmetic. (iii) $\Theta_s$ is a
blade. Against $E_{(k,a)}-E_{(k+1,a)}$ it gives the sign $(-1)^n$ for $a\ne s$ and $(-1)^{n-1}$ for $a=s$, so against $\tau_k$ the sign is $(-1)^{6n-1}=-1$.
The interior bivectors share no index with $\Theta_s$; the terms of $C^{kl}$ share none or two; a diagonal bivector $\Delta_{sb}$ shares one, so
$\Ad\Theta_s$ acts on the diagonal $\so{3,3}$ as the reflection of the direction $s$, an outer automorphism, which reverses $C_3$. (iv) $\tau_\sigma$
permutes the basis tensors of $S^{\otimes n}$ with the Koszul sign of $\omega$-parity, and $\omega$ has four even and four odd basis vectors. The image of
the symmetriser is spanned by orbit sums, and an orbit contributes iff no odd basis vector repeats, so its dimension is the number of
monomials of degree $n$ in four even and four odd variables, $\sum_k\binom{n-k+3}3\binom4k$; for the antisymmetriser exchange the roles of even
and odd, which gives the same number.
\end{proof}

Part (i) restates Proposition~7.3(iii) of \cite{RE} at $n=2$ and is fixed at every rung (e3-5a at $n=3,4$; e3-5e, relations and vacuum at
$n\le5$). Parts (ii)--(iv) are fixed (e3-5b; e3-5c, exact at $n\le4$, the $C_3$ eigenvalues in floating point to $10^{-6}$; e3-5e). Since $\tau_k$
and the natural action $\Lambda(s_k)$ of the transposition on $\Lambda(\C^n\otimes\C^3)$ implement the same automorphism of $\Cl(3n,3n)$, they differ by
a scalar, and the vacuum fixes it: $\tau_\sigma=\operatorname{sgn}(\sigma)\Lambda(\sigma)$. The Koszul swaps are therefore the permutation matrices of
Howe's $O(n)$, twisted by the sign character. Part (iii) generalises the mirror statement of \cref{e1:prop:stat}(ii) to every rung: the
Koszul sign does not sort the module by Lorentz type or by copy type, and its two sectors carry the same content up to the reflection of
one direction.

\begin{proposition}[The ungraded $n$-particle space]
\label{e3:prop:ungraded}
Let $P_\sigma$ permute the tensor factors of $S_n$ without signs, and let $W_n=\{J^{(1)}=\cdots=J^{(n)}\}$, of dimension $2\cdot4^n$, with $J:=J^{(1)}|_{W_n}$ as
imaginary unit, so that $W_n\cong S\otimes_\C\cdots\otimes_\C S=(\C^4)^{\otimes_\C n}$ (\cref{e1:prop:select}(iv)).
\begin{enumerate}[label=(\roman*),leftmargin=2.2em]
\item $P_\sigma J^{(k)}P_\sigma^{-1}=J^{(\sigma k)}$, so $P_\sigma$ preserves $W_n$ and is $\C$-linear there.
\item The totally antisymmetric part of $W_n$ under $P_\sigma$ is $\Lambda^n_\C\C^4$, of real dimension $2\binom4n=8,12,8,2,0$ at $n=1,\dots,5$, and the
totally symmetric part is $\mathrm{Sym}^n_\C\C^4$, of real dimension $2\binom{n+3}3=8,20,40,70,112$.
\item Under the diagonal $\lint$, $\Lambda^2_\C\C^4$ has $C_\Delta=0$ on $4$ and $3$ on $8$ dimensions, $2(0,0)\oplus(\half,\half)$; $\Lambda^3_\C\C^4$ has
$C_\Delta=\tfrac32$ on all $8$, a Dirac spinor; and $\Lambda^4_\C\C^4$ has $C_\Delta=0$, a scalar.
\end{enumerate}
\end{proposition}

\begin{proof}
(i) $P_\sigma(A_1\otimes\cdots\otimes A_n)P_\sigma^{-1}$ permutes the factors. (ii) Exact traces of $\frac1{n!}\sum_\sigma\operatorname{sgn}(\sigma)P_\sigma$ and of $\frac1{n!}\sum_\sigma P_\sigma$, composed with the projector
$2^{1-n}\prod_k(\id-J^{(k)}J^{(k+1)})$ onto $W_n$, as signed permutations. (iii) The projector onto the antisymmetric part, as an exact integer matrix,
satisfies the stated polynomial identities with $C_\Delta$. Since $\C^4=(\half,0)\oplus(0,\half)$,
$\Lambda^2=\Lambda^2(\half,0)\oplus(\half,0)\otimes(0,\half)\oplus\Lambda^2(0,\half)$, $\Lambda^3\C^4\cong(\C^4)^*\otimes\Lambda^4\C^4$ and $\Lambda^4\C^4=\det$.
\end{proof}

The representation theory is inherited \cite{WeinbergQFT}; its realisation on $W_n$ is fixed (e3-5d). With the plain exchange and
antisymmetry, rung $n$ on $W_n$ carries $n$ identical Dirac fermions in one momentum mode, and the Pauli principle holds: at most four, since a
Dirac spinor has four states, so there is no antisymmetric state at $n\ge5$ (inherited \cite{Pauli40}). The two exchanges therefore differ at
every rung, not only at $n=2$ (\cref{e1:prop:swap}). The Koszul exchange is the one the algebra supplies, since it permutes the vectors of
$\Cl(3n,3n)$. Its antisymmetric sector has dimension $88$ at $n=3$ and contains all three Lorentz types, and at $n=5$ it has dimension $360$
where the plain one is empty. It is not a statistics label. The plain exchange is a statistics label, but only on the constructed space $W_n$,
with the constructed generators $u_{(k,a)}$ as gamma matrices.

\subsection{The ladder table}
\label{e3:sec:table}

\begin{table}[t]
\centering
\footnotesize
\setlength{\tabcolsep}{4pt}
\begin{tabular}{@{}>{\raggedright\arraybackslash}p{4.7cm}ccccc>{\raggedright\arraybackslash}p{2.7cm}@{}}
\toprule
structure & $n=1$ & $2$ & $3$ & $4$ & $5$ & source\\
\midrule
$\dim V_n$, $\dim S_n$ & $6$, $8$ & $12$, $64$ & $18$, $512$ & $24$, $8^4$ & $30$, $8^5$ & ---\\
$\Cl(3n,3n)$ & \multicolumn{5}{c}{$M_{8^n}(\R)$ at every rung} & e3-1a ($n\le4$)\\
Fock modes & $3$ & $6$ & $9$ & $12$ & $15$ & e3-1c\\
spectator centraliser $\so{3n-1,3n-3}$ & $1$ & $28$ & $91$ & $190$ & $325$ & ES17 Thm~6.1; e3-6a\\
commutant of diagonal $\lint$, $\so n\oplus\so{2n}$ & $1$ & $7$ & $18$ & $34$ & $55$ & ES17 Thm~10.8; e3-6a\\
commutant of diagonal $\lG$, $\so{n,\C}$ & $0$ & $2$ & $6$ & $12$ & $20$ & ES17 Prop.~13.14; e3-6a\\
its compact part, $\so n$ & $0$ & $1$ & $3$ & $6$ & $10$ & e3-6a\\
rotor algebra $\uu n$ (in $\so{6n}$) & $1$ & $4$ & $9$ & $16$ & $25$ & ES17 Prop.~13.14; e3-6a\\
copy algebra, commutant of diagonal $\so{3,3}$ & $0$ & $1$ & $3$ & $6$ & $10$ & e3-4a\\
$C_{\rm cop}+C_{\rm diag}$ & $\tfrac{15}4$ & $9$ & $\tfrac{63}4$ & $24$ & $\tfrac{135}4$ & e3-4c\\
graded rest mass & $m$ & $\sqrt2\,m$ & $\sqrt3\,m$ & $2m$ & $\sqrt5\,m$ & e3-2a\\
ungraded rest energies $|m(n-2j)|$ & $m$ & $0,2m$ & $m,3m$ & $0,2m,4m$ & $m,3m,5m$ & e3-2b\\
$2\pi$ rotation, diagonal $\lint$ & $-\id$ & $+\id$ & $-\id$ & $+\id$ & $-\id$ & e3-3a\\
$2\pi$ rotation, diagonal $\lG$ & $+\id$ & $+\id$ & $+\id$ & $+\id$ & $+\id$ & e3-3b\\
copy-invariant four-vector & yes & no & yes & no & yes & ES17 Thm~13.15; e3-6b\\
selected imaginary unit (diagonal $\lint$) & yes & no & yes & no & yes & \cref{e1:prop:select}(iv)\\
invariant number of quanta (diagonal $\lG$) & no & yes & no & yes & no & \cref{e2:thm:invnumber}\\
invariant Grassmann mass kernel & $T_1$ & plane & $S_3$ & none & $T_5$ & \cref{e6:thm:ladder}(ii)\\
copy isospin & --- & integral & integral & $j_\pm\in\half\Z$ & --- & RE Prop.~9.6; e3-4d\\
projector $P_*$ & --- & --- & yes & --- & --- & ES17 Thm~8.1\\
$\Lambda^n_\C$ of the Dirac module on $W_n$ & Dirac & $2(0,0){\oplus}(\half,\half)$ & Dirac & $(0,0)$ & $0$ & e3-5d\\
Koszul-antisymmetric part of $S_n$ & $8$ & $32$ & $88$ & $192$ & $360$ & e3-5e\\
\bottomrule
\end{tabular}
\caption{The ladder. Dimensions of Lie algebras are real. ``Plane'' and ``none'': at $n\equiv2\pmod4$ the invariant Grassmann kernels form a plane,
at $n\equiv0$ there is none; $T_n$ and $S_n$ there are the temporal and spatial volumes of \cref{e6:thm:ladder}, not the spinor module. The copy
isospin is that of the six-direction copy algebra.}
\label{e3:tab:ladder}
\end{table}

\Cref{e3:tab:ladder} collects the rung-by-rung data. The Lie algebra rows are fixed (e3-6a, exact over $\Q$ at $n\le5$) and restate
Theorem~6.1, Theorem~10.8 and Proposition~13.14 of \cite{ES17}; the row of copy-invariant four-vectors restates Theorem~13.15 of \cite{ES17}, with
the case $n=3$ and the absence at $n=2$ re-verified (e3-6b). Read as a whole, the table says which structures single out which rungs.

\emph{No rung is singled out by the algebra.} $\Cl(3n,3n)=M_{8^n}(\R)$, its module is $S_n$, and every Lie algebra row is a polynomial in $n$. The
first structure with a period in $n$ is the symmetry of the invariant bilinear forms on $S_n$, of period four (\cref{e6:thm:ladder}(ii)). It is
Chevalley's pattern for spinors of $\R^{3n,3n}$, the dimension $6n$ running through the residues $6,4,2,0$ modulo $8$ with period four in $n$
(inherited \cite{Chevalley54}). Through it the $\so{3n,3n}$-invariant Grassmann mass kernel is unique exactly at odd $n$. That is the only
selection the algebra makes by itself, and it does not distinguish $n=3$ from the other odd rungs.

\emph{Odd rungs are singled out by the interior Lorentz action.} Under the diagonal $\lint$, odd $n$ is where the composite is a fermion
(\cref{e3:prop:spin}(i)), where an imaginary unit is selected (\cref{e1:prop:select}(iv)), where copy-invariant four-vectors and the Dirac block
exist (Theorem~13.15 and Proposition~10.10 of \cite{ES17}). \emph{Even rungs
are singled out by the binding action}: only there does a Lorentz-invariant number of quanta exist under the diagonal $\lG$
(\cref{e2:thm:invnumber}). \emph{The third rung is singled out by added projectors and by the low-rank coincidences of the diagonal algebra}: the
chiral Pati--Salam projector $P_*$ is constructed at $n=3$ (Theorem~8.1 of \cite{ES17}), the diagonal connection algebra there is
$\so3\oplus\so6\cong\su2\oplus\su4$ with colour $\su3$ inside $\so6$, and the copy isospin is integral (Proposition~10.10(iii) of \cite{ES17};
Proposition~9.6 of \cite{RE}). At $n=5$ the Dirac block carries one Standard-Model family per copy multiplet (Proposition~10.10(iv) of
\cite{ES17}). Apart from the period of the forms, every one of these selections comes from a Lorentz action, a projector or a gauge-algebra
identification added to the rung. None comes from the ladder.

\subsection{The step from \texorpdfstring{$n$}{n} to \texorpdfstring{$n+1$}{n+1}, and what going up the ladder means}
\label{e3:sec:answer}

\paragraph{What the step preserves.} By \cref{e3:prop:rung}(ii) the step keeps rung $n$ as the subalgebra $\Cl(3n,3n)\otimes\id$, with all of its
operators. Under the spectator action the interior four-plane stays in the first factor, so the tower is one Lorentz frame with nested
internal centralisers (Proposition~6.3 of \cite{ES17}). The diagonal module is tensored with $S$ (\cref{e3:prop:rung}(iv)), and the dual pair $(O(n),\so{3,3})$ sits inside
$(O(n+1),\so{3,3})$.

\paragraph{What the step adds.} Three Fock modes, $N_{n+1}=N_n\otimes\id+\id\otimes N_1$ (\cref{e3:prop:rung}(iii)). One frame: $V_{n+1}=V_n\oplus V$, which
under the diagonal $\lG$ is one more field strength $(1,0)\oplus(0,1)$ and under the diagonal $\lint$ one more potential $(\half,\half)$ with two
scalars (Section~13.5 of \cite{ES17}). One copy direction: the copy space $\R^n$ becomes $\R^{n+1}$, the copy algebra gains the $n$ rotations
$C^{k,n+1}$, and the Casimir constant grows by $\tfrac34(2n+5)$ (\cref{e3:thm:dual}(iii); e3-6c). The new constituent has its own commuting
copy, $(\Gamma_n\otimes\id)E_{(n+1,a)}$.

\paragraph{What the step changes.} The graded mass grows in quadrature and the ungraded energies add (\cref{e3:prop:mass}(iv)). The spin parity
under the diagonal $\lint$ flips, and that under the diagonal $\lG$ does not (\cref{e3:prop:spin}). Every parity-dependent selection of
\cref{e3:sec:table} switches on or off.

\paragraph{One particle or several.} Read together, the statements of this section answer the question in four sentences. (1) As given by
its vectors, every rung carries one particle: one Dirac mass $\sqrt n\,m$ along a unit vector of $\R^{3n,3n}$ with one clock, and copies
that are the components of a vector of the copy group $O(n)$. That group is the Howe dual of the diagonal $\so{3,3}$, its type is locked to
the $\so{3,3}$ type by $C_{\rm cop}+C_{\rm diag}=\tfrac34n(n+4)$, and its Lie algebra moves quanta between constituents without creating any. So
the copy index behaves as a flavour index of one particle, and going up the ladder adds a flavour direction and three modes to that particle
(fixed statements, read interpretively; \cref{e3:prop:mass,e3:thm:dual}). (2) The Lorentz content does not decide. The parity $(-1)^n$ under
the interior action is the parity of $n$ spin-$\half$ particles and equally that of one particle in an $n$-fold tensor product, and under
the field-strength action every rung is bosonic (fixed; inherited rule; \cref{e3:prop:spin}). (3) A rung carries $n$ particles in the sense of
ordinary quantum mechanics only with the structures of \cref{e1:sec:main}, now needed at every rung: the commuting copies
$u_{(k,a)}=\Gamma^{(1)}\cdots\Gamma^{(k-1)}E_{(k,a)}$ (fixed as an algebra), the additive Hamiltonian $H_u$ with its $n$ clocks, the common complex
space $W_n$ with $J|_{W_n}$, and the plain exchange with antisymmetry (constructed; the unit selected up to sign at odd $n$ only). With them
the rung carries $n$ identical Dirac fermions in one mode, with the Pauli bound $n\le4$ (fixed; inherited; \cref{e3:prop:ungraded}). (4) The
exchange the algebra supplies is the Koszul $\mathfrak S_n$, the permutation matrices of $O(n)$ up to sign, and it is not a statistics label: its
two sectors are mirror images and never empty (fixed; \cref{e3:prop:koszul}).

\paragraph{Which process goes up the ladder.} Within a rung, the copy algebra hops quanta among the $n$ constituents and conserves $N$
(\cref{e3:thm:dual}(ii); \cref{e2:prop:hop}). Between rungs there is no process in the algebra: rung $n+1$ is not an element or a state of
rung $n$, and the embedding $\Cl(3n,3n)\otimes\id$ makes rung $n$ the part of rung $n+1$ that does not see the last constituent. In the
two-particle reading the step is the tensor embedding $\psi\mapsto\psi\otimes\phi_0$, which needs a state $\phi_0$ that nothing selects
(constructed, as at $n=2$). Whether a dynamical mechanism selects $W_n$, the additive Hamiltonian and a rung is open.

\begin{table}[!htbp]
\centering
\footnotesize
\begin{tabular}{@{}>{\raggedright\arraybackslash}p{8.2cm}>{\raggedright\arraybackslash}p{4.5cm}>{\raggedright\arraybackslash}p{2.2cm}@{}}
\toprule
statement & class & checks\\
\midrule
\cref{e3:prop:rung}(i): $\Cl(3n,3n)=M_{8^n}(\R)$, volume $\omega^{\otimes n}$, $(-1)^N=(-1)^n\Gamma_n$ & fixed in the realisation; classification inherited & e3-1a\\
\cref{e3:prop:rung}(ii), (iii): step formulas, eight copies, commutant $\id\otimes M_8(\R)$; Fock modes & restated (RE Props.~7.1, 7.2; \cref{e1:prop:copies}) & e3-1b, e3-1c\\
\cref{e3:prop:rung}(iv): diagonal bivectors additive, $S_{n+1}=S_n\otimes S$ & fixed & e3-6c\\
\cref{e3:prop:mass}(i), (ii): $\pm\sqrt n\,m$; $m(n-2j)$ with $\binom nj4^n$; exact dressing & restated (\cref{e1:prop:rest}(vi), RE Prop.~8.11(iv)); multiplicities fixed & e3-2a, e3-2b\\
\cref{e3:prop:mass}(iii), (iv): one clock against $n$ clocks; the step & fixed & e3-2a, e3-2b, e3-6c\\
\cref{e3:prop:spin}(i), (iv): $e^{2\pi R_n}=(-1)^n$; $(-1)^{k_I}$; $C_\Delta$ spectra & restated (ES17 Props.~10.9, 10.10(i)); spectra fixed & e3-3a, e3-3c\\
\cref{e3:prop:spin}(ii), (iii): $e^{2\pi L_n}=+\id$; $(\half,\half)^{\otimes n}$; central element & fixed (extends ES17 Prop.~13.12) & e3-3b, e3-3c\\
\cref{e3:rem:copyparity}: spin parity against copy-invariant four-vectors & restated (ES17 Thms~10.6, 13.15); comparison interpretive & e3-6b\\
\cref{e3:thm:dual}(i): mutual commutants on vectors; $\eta$-weighting forced & fixed & e3-4a\\
\cref{e3:thm:dual}(ii): spin lifts, $\so n$, single-valued, Fock form & fixed; restated at $n=2,3$ (RE Props.~7.3, 9.6) & e3-4b, e3-4f\\
\cref{e3:thm:dual}(iii): $C_{\rm cop}+C_{\rm diag}=\tfrac34n(n+4)$; interior identity & fixed (constant new); interior restated (ES17 Thm~10.4) & e3-4c\\
\cref{e3:thm:dual}(iv): commutant at $n=2$ is the image of $\R[O(2)]$ & fixed & e3-4e\\
Howe's skew duality $(O(n),\so6)$ on $\Lambda(\C^n\otimes\C^3)$ & inherited \cite{Howe89} & e3-4f\\
\cref{e3:prop:joint}, \cref{e3:tab:joint}: joint decomposition at $n=2,3,4$ & fixed; $n=3$ copy spectrum restated (RE Prop.~9.6) & e3-4d\\
\cref{e3:prop:koszul}(i): Koszul $\mathfrak S_n$, Pin lift, commutation, vacuum sign & fixed; restated at $n=2$ (RE Prop.~7.3(iii)) & e3-5a, e3-5e\\
\cref{e3:prop:koszul}(ii)--(iv): $\mathfrak S_3$ content; mirror $\Theta_s$; sector dimensions & fixed & e3-5b, e3-5c, e3-5e\\
\cref{e3:prop:ungraded}: $W_n$, $\Lambda^n_\C\C^4$, $\mathrm{Sym}^n_\C\C^4$, Lorentz content & fixed in the realisation; inherited (exterior powers, Pauli) & e3-5d\\
\cref{e3:tab:ladder}: the ladder $n=1,\dots,5$ & fixed; restated (ES17 Thms~6.1, 8.1, 10.8, 13.15, Prop.~13.14; RE Prop.~9.6) & e3-6a, e3-6b, as listed\\
\cref{e3:sec:table}: no rung selected by the algebra; odd, even, third rung & interpretive (drawn from fixed statements); period four inherited & ---\\
\cref{e3:sec:answer}: step; one particle or several; process & fixed parts as cited; reading interpretive; tensor embedding constructed; mechanism open & e3-1b, e3-1c, e3-6c\\
\bottomrule
\end{tabular}
\caption{Claims of \cref{e3:sec:main}, their classes and their checks, which are \texttt{PASS} lines of \texttt{verify/e3\_ladder.py}.
``Restated'': fixed in \cite{ES17}, \cite{RE} or \cref{e1:sec:main}, re-verified here.}
\label{e3:tab:claims}
\end{table}
